% =============================================
% BYLAWS OF THE RENSSSELAER UNION EXECUTIVE BOARD
% !TEX root = main.tex
% !LW recipe=latexmk (lualatex)
% Compile with: latexmk main.tex
% =============================================
\documentclass{uniondoc}

\doctitle{The Rensselaer Union Executive Board Bylaws}
\orgname{Rensselaer Union Executive Board}
\orgshort{Executive Board}
\amendeddate{February 11, 2026}

\begin{document}

\article{Organization}

\begin{enumerate}
\item The name of this organization shall be the Rensselaer Union Executive Board, hereinafter referred to as the ``Executive Board'' or the ``Board.''
\item The Executive Board receives its authorities and powers from the Institute Board of Trustees.
\item The Executive Board shall issue decisions as well as implement and administer guidelines and procedures commensurate with its duties as defined in the Rensselaer Union Constitution and as specified in these Bylaws.
\end{enumerate}

\article{Membership}

\begin{enumerate}
\item A Representative shall be defined as any person who has voting rights in the Executive Board, as defined by the Rensselaer Union Constitution.
\item Executive Board membership shall consist of the President of the Union, Representatives, and appointed Officers.
\item Serving on the Executive Board shall be restricted to Activity Fee-paying members of the Union.
  \begin{enumerate}
  \item For this purpose, students not participating in the Arch summer semester shall be considered as Activity Fee-paying members of the Union from the conclusion of the spring semester until the commencement of the fall semester, provided that they met the requirement of Activity Fee-paying membership during the semesters before and after the summer.
  \end{enumerate}
\item No voting member of the Executive Board nor the President or Vice Presidents shall simultaneously hold a position on the Rensselaer Union Judicial Board, a voting position on the Rensselaer Union Student Senate, a voting position on the Rensselaer Union Undergraduate Council.
  \begin{enumerate}
  \item Furthermore, the President and Vice Presidents shall be prohibited from holding voting positions on the Undergraduate Class Councils and the Graduate Council.
  \item In the event of a conflict, including but not limited to after the confirmation of new members or election of a new President, the individual shall have 1 week to remove themselves from any conflicting positions.
  \end{enumerate}
\item The President shall be prohibited from serving as the president (or equivalent presiding officer position) or treasurer (or equivalent financial manager position) of any Union Funded or Recognized club or organization, except in acting capacities or were provided for in club governing documents for extenuating circumstances.
  \begin{enumerate}
  \item In the event of a conflict, including but not limited to after the election of a new President, the President shall have 1 week to remove themselves from any conflicting positions.
  \end{enumerate}
\end{enumerate}

\article{Officers}

\begin{enumerate}
\item The President of the Union, hereinafter referred to as the President, is the presiding officer of the Executive Board.
\item The Officers of the Executive Board shall be:
  \begin{enumerate}
  \item The chairs of any committee formed by these Bylaws or the Executive Board,
  \item The Vice President for Board Operations,
  \item The Vice President for Club Relations,
  \item The Vice President for Rules and Special Projects
  \item The Parliamentarian, and
  \item The Secretary.
  \end{enumerate}
\item All officers of the Executive Board, with the exception of the President and the chairs of any committee, shall be nominated by the President of the Union and confirmed by a two-thirds vote of the Executive Board.
  \begin{enumerate}
  \item The re-nomination of existing officers to previously-held positions shall still require the confirmation vote of the Executive Board.
  \end{enumerate}
\item All officers of the Executive Board, with the exception of the President and the chairs of any joint committee, shall serve at the pleasure of the President and may be removed from their positions at the discretion of the President or by two thirds vote of the Executive Board’s total voting membership.
\item All Vice Presidents shall:
  \begin{enumerate}
  \item Discharge the powers and duties of the President in the cases laid out in Article IX of these Bylaws or at the request of the President;
  \item Meet regularly with the President and the other Vice Presidents to discuss issues pertaining to the Rensselaer Union and the Executive Board; and
  \item Serve as a check on the President to ensure that the President faithfully fulfills and executes their duties.
  \end{enumerate}
\item The President shall nominate a Vice President for Board Operations from the Activity Fee paying membership of the Rensselaer Union. The Vice President for Board Operations shall:
  \begin{enumerate}
  \item Oversee the work of all Executive Board committees;
  \item Assist the President in the selection of committee chairs, and assume the duties of acting committee chairperson in the absence of a chairperson or identify an acting chairperson with the approval of the President;
  \item Assist the President in the organization of formal budgeting for the upcoming Fiscal Year;
  \item Coordinate annual documentation of committee work with the committee chairs to ensure continuity between years; and
  \item Meet with the chairs of all Executive Board committees and the President at least once per month to provide an update on committee work.
  \end{enumerate}
\item The President shall nominate a Vice President for Club Relations from the Activity Fee paying membership of the Rensselaer Union. The Vice President for Club Relations shall:
  \begin{enumerate}
  \item Oversee the club representative duties assigned to Executive Board Representatives and occasionally gain feedback on the fulfillment of these duties from club officers;
  \item Assist the President in the assignment of clubs for representation and assume the duties of assigned representative in the absence of an assigned representative if requested by the President;
  \item Organize and work with clubs placed on probation along with the Union staff and assigned Board representatives, if applicable;
  \item Ensure that all Union Clubs and Organizations have opportunity to contribute to any changes to the Rensselaer Union Guidelines \& Procedures;
  \item Coordinate annual documentation of specific club affairs with the respectively assigned representatives to ensure continuity between years; and
  \item Coordinate training sessions with Executive Board representatives and provide guidance to Board members where needed or requested.
  \end{enumerate}
\item The President shall nominate a Vice President for Rules and Special Projects from the Activity Fee-paying membership of the Rensselaer Union. The Vice President for Rules and Special Projects shall:
  \begin{enumerate}
  \item Maintain the Rensselaer Union Guidelines \& Procedures document, along with any supplementary appendices thereof;
  \item Serve as principal advisor to the President on these guidelines and procedures;
  \item Host working sessions open to Activity Fee-paying members of the Union to review and revise guidelines and procedures of the Executive Board or where instructed by the President or the Executive Board;
  \item Oversee the pursuit of special projects initiated by the President by collaborating with the appropriate committees, groups, students, and staff and developing recommendations to the President and the Executive Board where appropriate.
  \end{enumerate}
\item Vice Presidents may not simultaneously hold voting or other officer positions on the Rensselaer Union Executive Board.
\item The Vice Presidents may overrule certain decisions made by the President by a unanimous vote
  \begin{enumerate}
  \item There must be no fewer than two Vice Presidents in order for them to be able to overrule a decision of the President
  \item The decisions that may be overruled shall be established in the ``Override Regulation'', which shall be a separate document from the Executive Board Bylaws but shall be attached at the bottom for convenience.
    \begin{enumerate}
    \item These regulations will be adopted and may be modified by a 1/2 majority vote by the Board.
    \item The Executive Board shall have no regulations that enable the Vice Presidents to have override powers over the appointment of Representatives or Officers of the Executive Board.
    \item For each decision listed in the ``Override Regulations'', the source of this power must be cited from the relevant document. The President or the Vice President for Rules and Special Projects or designee may modify this policy to maintain the accuracy of wording should sections change.
    \item Should a decision power not be explicitly listed in the regulation, it shall not be overridable.
    \end{enumerate}
  \item For overrideable decisions made outside of a regular Board meeting, the President must inform the Vice Presidents of their intended decision, and the Vice Presidents have 16 hours to overrule their decision.
    \begin{enumerate}
    \item Should no decision be made within those 16 hours, the President may inform the relevant party of their decision.
    \item The Vice Presidents may waive their ability to overrule, in which case the President may inform the relevant parties.
    \end{enumerate}
  \end{enumerate}
\item The President shall nominate a Parliamentarian of the Executive Board from the Activity Fee paying membership of the Rensselaer Union. It shall be the responsibility of the Parliamentarian to:
  \begin{enumerate}
  \item Advise the presiding officer on matters of parliamentary procedure;
  \item Maintain an orderly meeting, abiding by rules of decorum
  \end{enumerate}
\item The President shall nominate a Secretary of the Executive Board from the Activity Fee paying membership of the Rensselaer Union. It shall be the responsibility of the Secretary to:
  \begin{enumerate}
  \item Record meeting attendance and minutes;
  \item Distribute minutes, attendance, committee reports, motions, and other printed material to the members of the Executive Board; and
  \item Ensure that minutes, attendance, public committee reports, and motions are made publicly available to any member of the Union within three business days after the date of such minutes, reports, or motions.
  \end{enumerate}
\item The President of the Union shall serve as the Treasurer of the Executive Board.
\item The President may, at their discretion, appoint other non-voting officers to assist with administering the business of the Executive Board.
  \begin{enumerate}
  \item All non-voting officers of the Executive Board shall serve at the leisure of the President of the Union.
  \item Any non-voting officer given comparable responsibilities to a Representative, such as representing assigned clubs, shall require the majority confirmation of the Student Senate before such duties can be conferred.
  \end{enumerate}
\item Upon the installment of a new President during Grand Marshal Week Elections, the Former President shall be regarded as a non-voting officer of the Executive Board from the installation of the new President until the conclusion of the current fiscal year.
  \begin{enumerate}
  \item The Former President shall only have the following assigned duties until the conclusion of the current fiscal year:
    \begin{enumerate}
    \item Assisting the President in transitioning into their role as President;
    \item Providing guidance on unfinished matters pertaining to the current fiscal year; and
    \item Attending Executive Board meetings as a guest for the remainder of the Fiscal Year to provide information and guidance where requested by the Executive Board, unless asked to not attend by the President.
    \end{enumerate}
  \item If the President is installed for a second term, the clauses pertaining to the Former President shall not apply for the commencement of the President’s second term.
  \end{enumerate}
\end{enumerate}

\article{Duties of Office}

\begin{enumerate}
\item Representatives shall either chair one Executive Board committee or sit on at least one Executive Board committee. The President and the Vice President for Board Operations shall determine if a Representative has fulfilled this requirement.
\item Each Representative and Officer shall be required to attend all Board meetings unless they obtain an excuse from the President.
\item Each Representative shall be assigned Union-funded organizations to represent to the Executive Board. The President and the Vice President for Club Relations shall determine if a Representative has fulfilled this requirement.
  \begin{enumerate}
  \item These assignments shall be determined by the President and may change at the discretion of the President.
  \item Representatives shall be expected to maintain frequent communication with the officers of their assigned clubs.
  \item Representatives shall be expected to thoroughly understand the goals and activities of their assigned clubs.
  \item In the absence of an assigned Representative for a club, the President and Vice President for Club Relations may discharge the responsibilities of that club’s representative.
  \end{enumerate}
\item Each Representative and Officer shall be expected to fulfill any additional requirements related to the business of the Board if such requirements are set by the President.
\item Each Representative and Officer must maintain a 2.5 GPA, but a Representative or Officer may appeal this requirement to the President. Academic standing will be verified with the Registrar at the time of application to the Board and at the commencement of each semester.
\item In handling matters pertaining to finances and operations, the Executive Board may be privy to confidential or sensitive information, designated as such by law, Institute policy, or at the discretion of the President. In these cases, each Representative and Officer shall be expected to uphold the confidentiality or privacy of the information.
\item Failure to meet any of the duties of office shall subject the offending Representative to removal procedures as outlined in Article IX, Section 3 of the Rensselaer Union Constitution.
\end{enumerate}

\article{Meetings}

\begin{enumerate}
\item The Executive Board shall meet at least once per month during the fall and spring semesters and with enough frequency and regularity to fulfill its duties to the Union.
\item The time and place of the meetings shall be determined by the President.
  \begin{enumerate}
  \item Meetings may be held virtually, through video conferencing software, if voting members are not able to attend in person as long as such a virtual meeting is publicized to members of the Union at least 48 hours in advance and members of the Union can join such a meeting if not closed.
  \end{enumerate}
\item Quorum shall consist of two-thirds of the voting membership of the Executive Board.
\item All meetings of the Executive Board shall be open publicly to the members of the Union, the Union Administration Staff, and any individuals invited by a member of the Executive Board.
  \begin{enumerate}
  \item The Executive Board may close a meeting with a two-thirds vote.
  \item Board can invite non-Board members to a closed meeting with a majority vote.
  \item All motions passed during the closed meeting must be made publicly available to members of the Union immediately upon request and published online according to standard procedure for publishing meeting minutes after the conclusion of the closed meeting.
  \item Minutes from a closed portion of a meeting may be released by a two-thirds vote of the Board.
  \end{enumerate}
\item The Executive Board may vote by a simple majority to open a closed meeting.
\item The Director of the Union shall provide a report during meetings of the Executive Board on the activities, facilities, and operations of the Union.
  \begin{enumerate}
  \item In the event of an inability to attend, the Director shall submit a report in writing to the President in advance of the meeting.
  \item In the event of a vacancy or extended absence of the Director, the President shall ask a member of the Union administrative staff to fulfill this responsibility.
  \item The President may remove any person who is disrupting a meeting.
    \begin{enumerate}
    \item The President’s decision may be overruled by the objection of any Representative.
    \end{enumerate}
  \item At any time, any member of the Executive Board may call for a vote to remove a person from a meeting. This requires a two-thirds vote of the Board and, if successful, cannot be overruled.
  \end{enumerate}
\item Emergency meetings may be called by the President, by unanimous decision of the Vice Presidents, or by one-quarter of the voting membership of the Executive Board.
  \begin{enumerate}
  \item If such a meeting is not called by the President, the Vice Presidents or requesting Representatives may name a Representative or Officer other than the President to preside over the meeting. The President must be notified of any meeting.
  \item The presiding officer is responsible for informing all Representatives and Officers at least twenty-four hours prior to the meeting as to the date, time, place, and purpose of the meeting. Normal procedures for running this meeting must be followed.
  \end{enumerate}
\item At the determination of the President, and with the consent of two-thirds of all Representatives, the Executive Board may adopt special rules of order for meetings. Written documentation of these special rules of order must exist and be available to all members of the Union. These rules of order shall be subordinate to these Bylaws. Any special rules shall expire at the end of the fiscal year in which they were established
\item The Executive Board may retain special rules of order into a new fiscal year by two-thirds vote.
\end{enumerate}

\article{Business}

\begin{enumerate}
\item In accordance with the Rensselaer Union Constitution, the Executive Board shall hear matters and make decisions pertaining to the finances, business affairs, operations, and facilities of the Rensselaer Union.
  \begin{enumerate}
  \item This collectively shall be known as the Jurisdiction of the Executive Board.
  \end{enumerate}
\item All motions brought before the Executive Board must be moved and seconded by voting members of the Executive Board, moved by the majority vote of the members of the Executive Board committee, or moved by a Vice President and seconded by a voting member of the Executive Board.
  \begin{enumerate}
  \item The author of the motion does not need to be the member moving the motion.
  \item A motion brought forward by an Executive Board committee does not require a second.
  \end{enumerate}
\item Motions pertaining to the Jurisdiction of the Executive Board must be submitted in writing.
  \begin{enumerate}
  \item The written motion must indicate the date of consideration by the Executive Board, the names and signatures of the individuals or the committee moving and seconding the motion, the outcome, and, if applicable, the vote count for the motion.
  \end{enumerate}
\item Motions pertaining to monetary transfers must clearly state the funds or organizational titles from which money is being taken and to which money is being transferred, along with the fiscal year in which the transfer is occurring.
\item Motions pertaining to purchases must indicate the vendor or company, the total sum of the purchase or an amount for the purchase not to exceed, and from what account the money will be appropriated.
  \begin{enumerate}
  \item The motion shall include a quote from the vendor for the purchase attached.
  \item The President may waive any of these requirements on a case-by-case basis at their discretion.
  \end{enumerate}
\item Motions pertaining to the business and finances of a Union club or organization originating from a request of that organization must include a written explanation from the organization’s officers, which should include the purpose of the purchase, why it is necessary, and the reasoning behind the chosen product or service.
  \begin{enumerate}
  \item The President may waive this requirement on a case-by-case basis at their discretion.
  \end{enumerate}
\item All motions shall be in accordance to any and all Union budgeting and monetary guidelines and procedures. Motions that contradict guidelines or procedures enacted by the Executive Board require at least a three-fifths majority vote of the Board to pass.
\item Motions pertaining to changes to guidelines or procedures within the Jurisdiction of the Executive Board must outline the exact changes proposed and shall require a two-thirds approval vote of the Executive Board to be enacted.
  \begin{enumerate}
  \item If the change is altering existing guidelines or procedures, that guideline or procedure and any applicable sections must be clearly identified. The original language that is being changed must be indicated, as well as the new language it is being changed to.
  \item If the change is establishing new guidelines or procedures, the new guideline or procedure must be attached to the motion being submitted.
  \item Approved changes to guidelines or procedures must be submitted to the Officers of the Union, as defined by the Rensselaer Union Constitution, and the Union Administration Staff within 24 hours of approval.
  \item Should the Grand Marshal determine that a revision exceeds the Jurisdiction of the Executive Board, they can decide within 48 hours of approval to defer the revision to the Student Senate for confirmation by a majority vote.
    \begin{enumerate}
    \item If such a determination is made, a statement noting the requirement of Student Senate confirmation to amend shall be included in the relevant sections.
    \end{enumerate}
  \item After the window of deference to the Student Senate has elapsed and after the confirmation of the Student Senate if deemed necessary, the President or the Vice President for Club Relations shall announce the changes to club officers and ensure the updated documents are made publicly available for members of the Union.
  \end{enumerate}
\item All guidelines and procedures enacted by the Executive Board shall be collated in a singular document entitled the Rensselaer Union Guidelines and Procedures and maintained by the Parliamentarian.
  \begin{enumerate}
  \item Supplementary or related documents, such as bylaws, tutorials, and reference materials, may be attached to this document in the form of appendices with the approval of the President of the Union, the unanimous consent of the Vice Presidents, or the majority vote of the Executive Board.
    \begin{enumerate}
    \item The Executive Board may remove an appendix by a majority vote.
    \item Each appendix shall clearly outline which position(s) or organization(s) is responsible for maintaining the information within that appendix, relevant contact information, and a last updated date.
    \item If the content of an appendix has its own amendment or modification procedure, the appendix shall be updated whenever this procedure occurs. Otherwise, appendices shall be updated with the approval of the Vice President of Rules and Special Projects or with the approval of the President.
    \end{enumerate}
  \end{enumerate}
\item The President may rule any motion out of order. The President’s decision may be overruled by a two-thirds vote of the Board.
  \begin{enumerate}
  \item If any two motions are similar in opinion or approach, the motion presented to the President first will be considered. The other will be ruled out of order.
  \item If any two motions present opposing opinions or approaches, the motion presented to the President first will be considered. If the first motion passes, the other will be ruled as out of order. If the first motion fails, the other will be heard.
  \end{enumerate}
\end{enumerate}

\article{Voting and Debate}

\begin{enumerate}
\item Only Representatives may vote on any business of the Executive Board. Unless otherwise specified, all votes will be determined by a simple majority vote.
\item The President may vote only in the case of a tie.
  \begin{enumerate}
  \item If a voting member of the Executive Board presides over a meeting, they cannot vote, except in the case of a tie, and do not count towards quorum.
  \item If the President (or presiding officer of a meeting) abstains, and the vote count remains tied, the vote shall not be considered a majority.
  \end{enumerate}
\item When hearing presentations, special orders, and motions pertaining to the Jurisdiction of the Executive Board, as defined in Article VI, a special procedure shall be utilized:
  \begin{enumerate}
  \item The President (or presiding officer of a meeting) shall introduce the proposal and presenter(s), if appropriate. The presenter(s) shall speak on the proposal, giving any relevant information and context.
  \item After the presentation, the President (or presiding officer of a meeting) shall open a discussion on the proposal. During this time, the Board may clarify its understanding of the proposal by questioning the presenter(s).
  \item When the discussion has been exhausted, if a motion on the topic had not been prepared in advance, the President (or presiding officer of a meeting) shall instruct the Representative responsible for the presenting organization to write a motion based on the discussion.
  \item Once one or more motions pertaining to the topic has been submitted to the President (or presiding officer of a meeting), all parties with a conflict of interest will be asked to leave the room.
    \begin{enumerate}
    \item A conflict of interest shall be defined by the Executive Board in the Rensselaer Union Guidelines \& Procedures document.
    \end{enumerate}
  \item The President (or presiding officer of a meeting) shall read the motion aloud and open the floor for discussion.
    \begin{enumerate}
    \item If a Representative suggests changes to the motion pertaining to its wording or language, but do not change its intent or merits, it may be amended by the author with the consent of the second.
    \item If a Representative suggests changes to the motion pertaining to its intent or merits, it may be amended by a majority vote of the Executive Board. This vote shall be conducted by a voice vote unless a show of hands or a roll call is requested by a Representative or the President.
    \item After any amendments are made, the President (or presiding officer of a meeting) shall read the amended motion aloud.
    \end{enumerate}
  \item Once all motions pertaining to the topic have concluded, parties with a conflict of interest shall be invited back into the room and the President (or presiding officer of a meeting) shall announce the outcome of all motions discussed.
  \end{enumerate}
\item Upon entering any discussion or debate, a queue will be established and maintained by the Vice President for Board Operations (or a designee in their absence selected by the presiding officer of a meeting).
  \begin{enumerate}
  \item The queue determines who has the floor to speak. The Vice President for Board Operations (or a designee in their absence selected by the presiding officer of a meeting) is responsible for acknowledging all those who wish to speak, and must announce each speaker before their speaking time begins.
  \item Anyone present in the room has the right to speak and may speak any number of times. However, those who have not spoken will be called upon before those who have spoken.
  \item By default, each speaking time shall be limited to 5 minutes.
  \item If the queue is on the topic of a presentation, designated presenters may respond to any questions directed at them. Speaking time limits still apply.
  \item Queues for any motion will proceed until the motion has left the floor. Other queues shall end when empty.
  \item A Representative may move to restrict the queue, requiring a second and majority approval.
    \begin{enumerate}
    \item By default, only Representatives, Officers, and Committee Chairs may be added to a restricted queue.
    \item Additional individuals to be given queue rights in the motion to restrict the queue or by majority vote of the Executive Board.
    \item If an individual without queue rights wishes to speak during a restricted queue, they will not be called upon until no members with queue rights remain on the queue.
    \end{enumerate}
  \item A Representative may move to close the queue, requiring a second and a two-thirds vote. No additional individuals may be added to a closed queue.
  \item A Representative may move to reopen a closed or restricted queue, requiring a second and a majority vote.
  \item When speaking time is limited, speakers can yield their time either to chair, forfeiting their time, or to another person in the room, transferring their remaining time.
    \begin{enumerate}
    \item A yielded speaker may decline to speak, forfeiting the remaining time.
    \end{enumerate}
  \item A yielded speaker may not yield time further.
  \item A Representative may move to call the question while they have the floor, requiring a second and a two-thirds vote. No additional debate is allowed and the motion in question will be brought to an immediate vote.
  \end{enumerate}
\item When voting on procedural or privileged motions, as defined by Robert’s Rules of Order, voting shall be conducted by voice, unless the President (or presiding officer of a meeting), a Vice President, or a Representative requests a show of hands or a roll call.
\item When voting on motions pertaining to the Jurisdiction of the Executive Board, as defined in Article VI, or ceremonial motions, voting shall be conducted by a show of hands, unless the President (or presiding officer of a meeting), a Vice President, or a Representative requests a roll call.
\item In a show of hands vote, the President (or presiding officer of a meeting) will request the number in favor, then the number opposed, then the number of abstentions. All Representatives present must vote or indicate their decision to abstain.
  \begin{enumerate}
  \item Once all votes have been tallied, the President (or presiding officer of a meeting) will announce the result of the vote.
  \item In the case of a tie, the President (or presiding officer of a meeting) shall announce their vote, or their decision to abstain, and then the outcome of the vote.
  \item The vote count on any motion shall be recorded on the motion, and in the minutes by the Secretary, and it shall be made publicly available to the members of the Union with the motion.
  \end{enumerate}
\item If a roll call is requested, the President (or presiding officer of a meeting) shall read the names of Representatives ordered alphabetically by surname.
  \begin{enumerate}
  \item Each Representative shall vote or pass, in which case their name will be read again after the first completion of the roll call. At this time, the Representative must vote or abstain. Refusal to vote shall be determined an abstention.
  \item The result of the roll call shall not be determined until all Representatives present have voted or abstained.
  \item The names of those in favor and those against shall be recorded on the motion and in the minutes by the Secretary and made publicly available to the members of the Union with the motion.
  \end{enumerate}
\item In cases where the Executive Board must conduct business, but is unable to meet, the President may entertain a motion to be voted on by means of an electronic vote through a proxy meeting.
  \begin{enumerate}
  \item An electronic vote shall not be held if one quarter of the Executive Board Representatives object.
  \item A Representative shall make a motion via email sent to every Representative and the President.
    \begin{enumerate}
    \item The President may use another medium for an electronic vote with the approval of the Vice Presidents as long as all Representatives and Officers of the Executive Board are included in that medium.
    \end{enumerate}
  \item Electronic motions do not require a second.
  \item Once the electronic motion has been submitted via email, the President shall establish a voting period of no less than 24 hours and no greater than 7 days.
    \begin{enumerate}
    \item Once this time period has been set, it cannot be altered.
    \item A quorum of the Executive Board must be reached in the set time period for the electronic vote to be valid and binding.
    \end{enumerate}
  \item Representatives shall submit their vote by responding to the President and all Executive Board Representatives and Officers with a message that clearly indicates an affirmative vote, a negative vote, or an abstention.
  \item Upon the conclusion of the set voting period, the President shall respond with the vote count, whether a quorum was reached, and the outcome of the vote.
  \item If passed, the motion shall be made publicly available to the members of the Union immediately following the conclusion of the set voting period.
  \item The discussion thread of any electronic medium where the motion in question was discussed shall constitute the minutes of the electronic vote. These discussion threads shall be made publicly available to the members of the Union as minutes by the Secretary with the omission of any private or personal information of the participants within 24 hours of the expiration of the set voting period.
  \end{enumerate}
\end{enumerate}

\article{Committees}

\begin{enumerate}
\item The Executive Board shall have two types of committees: standing and special.
  \begin{enumerate}
  \item Standing committees must be defined in these Bylaws with defined purposes and area of jurisdiction, and they shall persist across fiscal years.
  \item Special committees, also known as temporary, ad-hoc, or select committees, may be created by majority vote of the Executive Board.
    \begin{enumerate}
    \item These committees shall be given a specific purpose and have such powers granted to it by the Executive Board, to support its given purpose. These powers shall not exceed the powers of the Executive Board.
    \item The Board may impose a time restriction on a special committee. If no such restriction is established then the committee shall be disbanded upon the completion of their established goal or the end of the fiscal year in which the committee was formed, whichever occurs first.
    \item The Executive Board may extend the tenure of a special committee by majority vote, and the Executive Board may dissolve a special committee at any time by two-thirds vote.
    \end{enumerate}
  \end{enumerate}
\item The Student Senate and the Executive Board may designate their respective committees as joint committees.
\item Standing joint committees shall be defined both in these Bylaws and in the Bylaws of the Rensselaer Union Student Senate. The Executive Board shall have the following standing committees:
  \begin{enumerate}
  \item The Business Operations Committee (Biz Ops);
  \item The Club Operations Committee (Club Ops);
  \item The Marketing and Strategy Committee (M\&S);
  \item The Multicultural Leadership Council (MLC);
  \item The Union Programs and Activities Committee (UPAC); and
  \item The Executive Communications Committee (ExComm).
  \end{enumerate}
\item The Executive Board shall, with the Student Senate, have the following joint standing committee:
  \begin{enumerate}
  \item The Union Annual Report Committee.
  \end{enumerate}
\item Each committee shall be chaired by a member of the Rensselaer Union.
  \begin{enumerate}
  \item The chairperson of each non-joint committee, with the exception of the chairperson of the Union Programs and Activities Committee, shall be nominated by the President and confirmed by two-thirds vote of the Executive Board.
  \item The chairperson of each joint committee, with the exception of the chairperson of the Union Annual Report Committee, shall be nominated by the Grand Marshal and the President of the Union and confirmed by two-thirds vote of both the Student Senate and the Executive Board.
  \item The chairpersons of each non-joint committee shall serve at the pleasure of the President.
  \item The chairpersons of each joint committee shall serve at the pleasure of the Grand Marshal and the President.
  \end{enumerate}
\item A vice chairperson for each committee shall be selected by the committee chairperson.
  \begin{enumerate}
  \item The vice chairperson shall be appointed within four committee meetings after the committee chair’s appointment.
  \item The vice chairperson of each committee shall be responsible for chairing meetings of the committee when the chairperson is unable to do so.
  \item Vice chairpersons shall serve at the pleasure of the committee chairperson and of the Vice President for Board Operations.
  \item Vice chairpersons of joint committees shall serve at the pleasure of the committee chairperson, of the Vice President for Board Operations, and of the Grand Marshal.
  \end{enumerate}
\item Committee membership shall be open to all members of the Rensselaer Union.
  \begin{enumerate}
  \item Any committee member who is not a member of the Executive Board must be approved by the committee chairperson before receiving membership in the committee.
  \item Members of a committee shall serve at the pleasure of the committee chairperson and of the Vice President for Board Operations and may be removed at any time for good cause.
    \begin{enumerate}
    \item Good cause for removal includes, but is not limited to, neglecting the duties of membership, failure to adhere to the rules of order during meetings, or failure to adhere to these Bylaws when conducting business in the name of the committee.
    \item Members are expected to dutifully attend meetings of the committee and complete committee assignments within reasonable deadlines.
    \end{enumerate}
  \end{enumerate}
\item For each project undertaken by a committee, a project lead shall be selected by the chairperson of the committee. Project leads shall serve at the pleasure of the chairperson and of the Vice President for Board Operations.
\item The Vice President for Board Operations shall maintain a list of committee chairpersons, vice-chairpersons, and project leads for all committees.
\item During meetings, each committee shall follow rules of order determined at the discretion of its Chair. Such rules of order shall be communicated to the President and Vice Presidents and made publicly available to the members of the Union immediately following their adoption and before any official committee business may commence.
  \begin{enumerate}
  \item The Secretary shall ensure that any such rules of order are made publicly available to members of the Union.
  \end{enumerate}
\item All committees shall be empowered to create subcommittees to aid them in the performance of their functions. The President and Vice President for Board Operations must be consulted before the creation of any subcommittee. The Committee Chair shall be an ex-officio member of any such subcommittee.
\item Each committee shall meet at least twice per month during the fall and spring semesters or with enough frequency to satisfy its obligations to the Executive Board.
  \begin{enumerate}
  \item Attendance shall be recorded at all committee meetings and submitted to the Vice President for Board Operations.
  \end{enumerate}
\item Each committee chair shall be held accountable for their committee’s performance.
\item The Club Operations Committee shall:
  \begin{enumerate}
  \item Represent the interests of non-funded Union clubs and organizations to the Executive Board when needed;
  \item Review all applications for the classification of new organizations and changes in classifications for old ones and make recommendations to the Board;
  \item Review and approve the constitutions of all clubs operating within the Rensselaer Union;
  \item Review any application by a club for funding in accordance with the Rensselaer Union Guidelines \& Procedures and administer the club’s starter budget, pending approval of the Executive Board;
  \item Receive updates from the Vice President for Club Relations on club representational matters; and
  \item Review and hold clubs accountable in accordance with the Rensselaer Union Guidelines and Procedures.
  \end{enumerate}
\item The Business Operations Committee shall:
  \begin{enumerate}
  \item Assess the state of all retail services under the direction of the Executive Board, develop ways to improve such services with the objective of improving the quality and diversity of goods and services available to the Rensselaer Community, and report its findings and recommendations to the Executive Board;
  \item Investigate new sources of revenue to the Union, maintain all existing sources of revenue excluding the Activity Fee, and report its findings and recommendations to the Executive Board;
  \item Investigate the current state of all Union facilities and report its findings and recommendations to the Executive Board;
  \item Oversee the facilities of the Rensselaer Union and recommend allocation of space within the Union to the Executive Board; and
  \item Oversee the pursuit of special projects initiated by the President by collaborating with the appropriate committees, groups, students, and staff and developing recommendations to the President and the Executive Board where appropriate.
  \end{enumerate}
\item The Marketing \& Strategy Committee shall:
  \begin{enumerate}
  \item Oversee procedures pertaining to marketing, communication, and social media for the Union;
  \item Consider matters pertaining to future operations of the Union and developing strategies to guide the Union in achieving those futures;
  \item Guide Union organizations and staff in promoting the Union and related activities, services, and events; and
  \item Maintain a cohesive Union brand.
  \end{enumerate}
\item The Multicultural Leadership Council shall:
  \begin{enumerate}
  \item Organize events on and promote awareness of, for the benefit of the entire student body, the eight pillars of diversity and identity: ability, age, ethnicity, gender, race, religion, sexual orientation, and socioeconomic status;
  \item Foster and encourage forums of discussion where issues of diversity and inclusion can be exposed, reflected upon, and actively analyzed;
  \item Oversee the affairs of the Multicultural Lounge, in conjunction with the Business Operations Committee, working to maintain and improve the space in consideration of students’ needs and desires for the space; Coordinate business and relationships between identity-based organizations recognized by or affiliated with the Union to further their collective interests and promote inter organization collaboration;
  \item Lead the operation of the Multicultural Round Table, which shall:
    \begin{enumerate}
    \item Consist of the presiding officer, or their designee, of each identity-based organization recognized by or affiliated with the Union,
    \item Meet monthly during the fall and spring semesters to discuss affairs, interests, priorities, and updates from identity-based organizations, and
    \item Operate in any additional capacities outlined in the Council’s Bylaws;
    \end{enumerate}
  \item Establish subcommittees, at the Council’s discretion, to assist in the fulfillment of the Council’s purposes:
    \begin{enumerate}
    \item Each subcommittee shall be chaired by a member of the Council selected by the chairperson of the Council and approved by vote of the Council and the President of the Union;
    \end{enumerate}
  \item Consist of a set membership of up to 8 voting members with a minimum of 6 voting members and non-voting membership open to any member of the Union:
    \begin{enumerate}
    \item Voting members will be required to maintain at least a 2.5 GPA for their entire term on the Council.
    \item The voting members shall be selected by the previous voting membership within two weeks of the conclusion of Grand Marshal Week, in consultation with the newly-elected President and the Union staff, and shall be approved by vote of the Executive Board,
    \item Terms for voting members shall begin upon the approval of the Executive Board and shall serve until replaced,
    \item Should a voting member position become vacant, the Council’s chairperson shall appoint, in consultation with the President and the Union staff, a member of the Union to fill the vacancy with the approval of the Executive Board, and
    \item Should the council be unable to select new voting members, the President shall work with the Council’s chairperson, if filled, and the Union staff to select new voting members for Executive Board approval;
    \end{enumerate}
  \item Be chaired by a member selected by majority vote of the Council’s voting members and approved by the President and a two-thirds majority vote of the Executive Board:
    \begin{enumerate}
    \item Should the chairpersonship remain vacant and the Council is unable to select a chair, the President may appoint any member of the Union to the chairpersonship; and
    \item Maintain bylaws detailing operations and the roles of any subcommittees that are subsidiary to these Bylaws and are approved by the Council and the Executive Board.
    \end{enumerate}
  \end{enumerate}
\item The Union Programs and Activities Committee shall:
  \begin{enumerate}
  \item Organize events for the benefit of the entire student body, specifically by:
    \begin{enumerate}
    \item Creating such events of its own origin,
    \item Providing support for clubs and organizations hosting such events, and
    \item Possessing and operating resources necessary to make such events possible;
    \end{enumerate}
  \item Establish subcommittees to better maintain and utilize specific facilities and resources with regards to their programming activities.
    \begin{enumerate}
    \item Each subcommittee shall be chaired by a person or persons elected from its membership, and
    \item The chairperson(s) of each subcommittee or their duly appointed representatives shall be an officer of the committee; and
    \end{enumerate}
  \item Maintain the following standing subcommittees:
    \begin{enumerate}
    \item Cinema,
    \item Comedy,
    \item Concerts,
    \item Grand Marshal Week, and
    \item Winter Carnival.
    \end{enumerate}
  \item Be chaired by a person selected by majority vote of its members and approved by the President and a two-thirds majority vote of the Executive Board. If the chairpersonship remains vacant and the Committee is unable to select a chair, the President may appoint any member of the Union to the chairpersonship.
  \item Maintain bylaws that are approved by the committee and the Executive Board detailing operations and subcommittee roles.
  \end{enumerate}
\item The Executive Communications Committee, hereby referred to as ExComm, shall be responsible for promoting the initiatives and activities of the Executive Board.
  \begin{enumerate}
  \item Its chairperson may act as the spokesperson for the Executive Board at the discretion of the President of the Union, responsible for press releases and other communication with media outlets.
  \item It shall maintain bylaws that are approved by the committee and the Executive Board, detailing operations and subcommittee roles.
  \item It shall make use of a variety of communication media to publicize the initiatives and activities of the Executive Board.
    \begin{enumerate}
    \item To this end, it shall coordinate with the Executive Board’s standing committees to promote current projects and initiatives.
    \item It shall provide the student body with accurate information on the activities, initiatives, policies, and projects of the Executive Board.
    \end{enumerate}
  \item It shall make efforts to ascertain the concerns and opinions of students regarding matters under discussion in the Executive Board.
    \begin{enumerate}
    \item It may solicit student concerns and opinions through public relations initiatives.
    \item It shall inform other Committees of student concerns pertaining to the purview of each Committee.
    \end{enumerate}
  \item It shall foster, among the student body, informed discussion of issues currently facing the Executive Board.
    \begin{enumerate}
    \item It shall, when necessary, create and distribute informational media in conjunction with its ongoing public relations initiatives.
    \item It shall create and maintain platforms for student inquiry into the activities of the Committees.
    \item It may collaborate with the Marketing and Strategy Committee (M\&S) when necessary.
    \end{enumerate}
  \item It shall promote cohesion between the Executive Board and other Union bodies.
    \begin{enumerate}
    \item It shall be the responsibility of the chairpersons, or designated representative, of each Executive Board committee to correspond with both the ExComm chairperson and the President of the Union on all publicity information.
    \item The Director of Public Relations of the Undergraduate Council, a designee of the Graduate Council, and the chairperson of the Judicial Board may correspond with ExComm and its chairperson if they wish to use the resources and platforms of the committee to promote the information and initiatives of their respective bodies.
    \item The Chair of ExComm and their equivalents on the Student Senate, Judicial Board, Undergraduate Council, and Graduate Council, or a designee if there is no such equivalent, shall meet at least monthly during the academic year to ensure cohesion between different Union bodies in their communications and messaging.
    \end{enumerate}
  \end{enumerate}
\item The presiding officer, with a 1/2 vote of the Board, may resolve the Board into an Executive Committee for Budgeting and Hiring during a regular meeting.
  \begin{enumerate}
  \item The chair for this committee shall be the President of the Union, or another designated member to be approved by a 1/2 vote of the Board.
  \item Upon resolving into this committee, the room shall be closed to all except the membership of this committee.
  \item The membership of this committee shall be the membership of the Board and Union Administrators.
    \begin{enumerate}
    \item Only voting members may vote on motions proposed in this committee.
    \item Non-voting members of the Board may move and second motions while in this committee. They may vote to convey their opinion, but may not vote on final motions.
    \item Union Administrators may not move or second motions.
    \end{enumerate}
  \item The purpose of this committee shall be to expedite the budgeting and hiring process of the Board.
  \item The votes and decisions of this committee are to be considered the votes and decisions of the Board.
  \item While in this committee, the Board shall adopt special rules of procedure.
    \begin{enumerate}
    \item The Board will follow procedure as if in a committee.
    \item The chair, if they so desire, may instruct the Parliamentarian to maintain a queue.
    \end{enumerate}
  \item Upon the conclusion of the budgeting session, or by a majority vote of the committee, the committee shall rise, and normal Board procedures shall resume.
  \end{enumerate}
\item The Union Annual Report Committee shall fulfill duties and responsibilities outlined in the Bylaws of the Rensselaer Union Student Senate.
\item Positions on any Student Senate, Institute, or otherwise non-Executive Board committees reserved for Executive Board representation shall be filled by members of the Executive Board appointed by the President and confirmed by majority vote of the Executive Board.
\end{enumerate}

\end{document}
